\scanpage{72}
% archive sheet number: 71

\uline{Th} \textcircled{I} Soit $X$ une courbe analytique régulière, $\deleted{\unsure{\mathcal P_X|\mathcal E_X}}\inserted{\unsure{\mathcal E|\mathscr E}_X}\ (\unsure{\mathcal P|\mathscr P}_X)$ le faisceau des applications étales d'ouverts de $X$ dans $\mathbb E^1(\mathbb C)$ $(\mathbb P^1(\mathbb C))$. Il existe un morphisme

\[
G:\mathcal E_X\times\mathcal E_X\longrightarrow\underline\omega_{X/\mathbb C}^{\otimes2}
\qquad(\underline\omega_{X/\mathbb C}=\underline\Omega^1_{X/\mathbb C})
\]

ayant les propriétés suivantes

a) $G$ est un op. différentiel (pas néc linéaire), i.e. en termes d'une uniformisante locale $t$ \inserted{(déf. sur un ouvert $U$)} on a

\[
G(f,g)=R(f,f',\ldots,f^{(n)},g,g',\ldots,g^{(n)})\,dt^2
\]

$R$ est une fonction rationnelle en $\unsure{2n|2m}$ variables, à coefficients dans $\Gamma(\unsure{U|X},\mathcal O_U)$. De plus, \deleted{\unsure{le dénominateur|le dénominateur n'est}} \inserted{si $R=P/Q$, $P$ et $Q$ polynômes,} $Q$ non identiquemt nul, \deleted{si on substitue} ce qui implique \deleted{\unsure{est tel qu'il n'a pas de pôles zéros}} que $Q(f,\unsure{g'|f'},\ldots,f^{(n)};g,g',\ldots,g^{(n)})$ n'est pas identiquemt nul, \deleted{\unsure{quel que soit le choix de $f$}} \inserted{\unsure{$f$ et $g$}} on suppose que $Q$ ne s'annule identiquemt pour aucuns $f$ et $g$ avec $f$ et $g$ non loc. constantes ni l'une ni l'autre, et que $P(f,\ldots,g,\ldots)/Q(f,\ldots,g,\ldots)$ n'a \deleted{pas} de pôles qu'aux pts où $f',g'$ ont des zéros.

[On doit pouvoir montrer que, si on prend $P,Q$ premiers entre eux, on \deleted{aura} \inserted{peut prendre} $Q={f'}^{p}{g'}^{q}$, donc

\[
G(f,g)=\frac{1}{{f'}^{p}{g'}^{q}}\,P(f,f',\ldots,f^{(n)};g,g',\ldots,g^{(n)})\quad]
\]

\[
\text{b)}\quad G(f,g)+G(g,h)=G(f,h)
\]

\[
\begin{aligned}
\text{c)}\quad G(f,ag)&=G(f,g)\deleted{\unsure{{}+G(g,\tfrac1g)|{}=G(f,\tfrac1g)}}&
 a\in\mathbb C^*\quad\deleted{\text{si }a\in\mathbb C^*,b\in\mathbb C}\\
G(f,g+b)&=G(f,g)& b\in\mathbb C\quad\deleted{\unsure{\text{et }g|\text{et}}}\\
G(f,1/g)&=G(f,g)&\text{si $g$ sans zéros}
\end{aligned}
\]

d) \deleted{\unsure{Pour $f$ fixé,|Deux $f$ fixés,}} Si $G(f,g)=0$ et $f,g$ coïncident à l'ordre 2 en un point $x\in X$ (\unsure{i.e.} $f,g\in\Gamma(U,\underline{\mathcal O}_U)$) alors $f=g$ si $U$ connexe.

% margin note written sideways
\margin{e) Si $f\in\Gamma(U,\mathcal E_U)$, le morphisme $g\longmapsto G(f,g)$ de $\mathcal E_U$ dans $\underline\omega_U^{\otimes2}$ est un épi}

% margin note written sideways
\margin{\uline{NB} c) implique que $G$ \unsure{soit défini|se définit|se factorise en} sur $\mathcal P_X\times\mathcal P_X\longrightarrow\underline\omega_X^{\otimes2}$, \deleted{\unsure{et se factorise en}} \inserted{\unsure{et signifie qu'il est}} $\,{}_H\backslash\mathcal P_X\times {}_H\backslash\mathcal P_X\longrightarrow\underline\omega_X^{\otimes2}$, où $H=\mathrm{GP}(1,\mathbb C)$. b) donne la relation analogue pour $\,{}_H\backslash\mathcal P_X$, d'où torseur $\mathcal T_X$ et $\,{}_H\backslash\mathcal P_X\xrightarrow{c}\mathcal T$ avec $G(f,g)\unsure{\equiv|=}c(f)-c(g)$; Alors d) signifie que $\unsure{{}_H\backslash\mathcal P_X|\mathcal P_X}\xrightarrow{c}\mathcal T_X$ est mono, et e) que c'est épi}

Deux tels morphismes $G'$ et $G$ sont liés par

\[
G'=\varphi G,\qquad\text{avec }\varphi\in\Gamma(X,\underline{\mathcal O}_X^*)
\]

\textcircled{II} Si, au lieu de prendre $X$ fixe, on prend $X$ variable, avec morphismes étales entre des $X$, on a $=$ résultat, et l'indétermination sera

% margin note written sideways
\margin{Ici, au lieu des conditions d), e) il suffit de supposer d) \deleted{\unsure{que $G$ est fidèle}}}

\[
G'=c\,G\qquad c\in\mathbb C^*
\]

\unsure{TSVP|TSVo}
